\scanpage{225}
% slanted
\margin{Feuille volante paumée ...}

\uline{Prop 6} Soit $\unsure{\mathcal B|B|\mathfrak B}$ une catégorie $k$-lin karoubienne. Alors le foncteur
\[
 ND_\bullet\colon\mathcal B_*^{\deleted{\unsure{*|\widehat{\ }}}}\longrightarrow\mathcal B_\bullet
\]
est donné (à isom can. près) par
\[
 ND_i(L_*)\simeq\operatorname{Hom}_{k,*}(\underbrace{\Lambda^i\unsure{\Phi|\mathscr L|\mathfrak L}_*}_{\Phi^i_\bullet},L_*)
\]
(pour la définition du $\operatorname{Hom}_{k,*}$, cf pag. 3, 5), \deleted{l'opérateur différentiel étant \unsure{les|la}} 
\[
 ND_\bullet(L_*)\simeq\operatorname{Hom}_{k,*}(\underbrace{\Lambda^\bullet\Phi_*}_{\Phi^\bullet_*},L_*)
\]
l'opérateur diff. sur $ND_\bullet(L_*)$ provient de l'op. différentiel sur \unsure{$T_*1$|Tout|$T_*$ et} sur $\Lambda^\bullet\Phi_*$.

En effet, \deleted{\unsure{comme}} \deleted{\unsure{on a pour}} soit $L_\bullet=ND_\bullet(L_*)$, et \deleted{\unsure{soit}} $\Phi(i)_\bullet=ND_\bullet(\deleted{\Phi}\Lambda^i\Phi_*)$, alors on a par \inserted{\unsure{la corresp.|l'équiv.|la symétrie} de \unsure{Dold, Puppe|Dold-Puppe}}
\[
 \operatorname{Hom}_{k,*}(\Lambda^i\Phi_*,L_*)\simeq\operatorname{Hom}_{k,\bullet}(\Phi(i)_\bullet,L_\bullet)
\]
Or $\Phi(i)_\bullet=C_\bullet^{i-1}$ pour $i\geq1$ (prop 4), et on \inserted{$\Phi(0)_\bullet=k$ placé en degré 0} a bien
\[
 \left\{\begin{aligned}
 \operatorname{Hom}_{k,\bullet}(C_\bullet^{i-1},L_\bullet)&\simeq L_i\quad(i\geq1)\\
 \operatorname{Hom}_{k,\bullet}(\unsure{(}k[0]_\bullet,L_\bullet)&\simeq L_0
 \end{aligned}\right.
\]
et il reste à faire de petites vérifications de compatibilité \unsure{...}.

\uline{Cor} Le foncteur $ND^\bullet\colon\mathcal B^*\to\mathcal B^\bullet$ s'explicite comme
\[
 ND^\bullet(L^*)\simeq\Lambda^\bullet\Phi_*\otimes_{k,*}L^*.
\]
[où pour un $M_*\in\mathrm{Ab}_{k*}$ et $L^*\in\mathcal B^*$,
$M_*\otimes_{k*}L^*$ désigne l'objet de $\uline{\operatorname{Hom}}(\mathcal B,(\unsure{Abgr|\mathrm{Ab}}))$
\[
 b\longmapsto\operatorname{Hom}_{\mathrm{Ab}_{k*}}(M_*,\operatorname{Hom}_{\mathcal B}(L^*,b))]
\]
\margin{$\big|$}

Il suffit d'appliquer la prop. 6 à $\mathcal B^\circ$.

% archive sheet number: 224

